\usepackage[T1]{fontenc}
\usepackage{lmodern}
\input{glyphtounicode}
\pdfgentounicode=1
\pdfglyphtounicode{parenleftbig}{0028}
\pdfglyphtounicode{parenrightbig}{0029}
\pdfglyphtounicode{bracketleftbig}{005B}
\pdfglyphtounicode{bracketrightbig}{005D}
\pdfglyphtounicode{vextendsingle}{23D0}
\pdfglyphtounicode{vextenddouble}{2016}
\pdfglyphtounicode{parenleftBig}{0028}
\pdfglyphtounicode{parenrightBig}{0029}
\pdfglyphtounicode{parenleftbigg}{0028}
\pdfglyphtounicode{parenrightbigg}{0029}
\pdfglyphtounicode{braceleftbigg}{007B}
\pdfglyphtounicode{bracerightbigg}{007D}
\pdfglyphtounicode{parenleftBigg}{0028}
\pdfglyphtounicode{parenrightBigg}{0029}
\pdfglyphtounicode{bracketleftBigg}{005B}
\pdfglyphtounicode{bracketrightBigg}{005D}
\pdfglyphtounicode{summationtext}{2211}
\pdfglyphtounicode{integraltext}{222B}
\pdfglyphtounicode{summationdisplay}{2211}
\pdfglyphtounicode{integraldisplay}{222B}
\pdfglyphtounicode{hatwide}{02C6}
\pdfglyphtounicode{tildewide}{02DC}
\pdfglyphtounicode{bracketleftBig}{005B}
\pdfglyphtounicode{bracketrightBig}{005D}
\pdfglyphtounicode{radicalbig}{221A}
\pdfglyphtounicode{radicalBig}{221A}
\pdfglyphtounicode{radicalBigg}{221A}
\pdfglyphtounicode{mapsto}{007C}
\pdfglyphtounicode{arrowhookleft}{21AA}

\usepackage[margin=1in]{geometry}
\usepackage{amsmath,amssymb,amsthm,mathtools}
\usepackage{microtype}
\usepackage{enumitem}
\usepackage{booktabs,longtable,array}
\usepackage{xcolor}
\usepackage{tikz}
\usetikzlibrary{arrows.meta,positioning,fit}
\usepackage[numbers,sort&compress]{natbib}
\definecolor{linkblue}{RGB}{28,64,104}
\usepackage[colorlinks=true,linkcolor=linkblue,citecolor=linkblue,urlcolor=linkblue,breaklinks=true]{hyperref}
\usepackage{bookmark}
\hypersetup{pdftitle={A Charged Reduction of the Spacetime Penrose Inequality in Spatial Dimensions at Least Four},pdfsubject={Purely electric initial data, neutral Penrose input, and original-data rigidity},pdfauthor={OpenAI}}
\numberwithin{equation}{section}
\newtheorem{theorem}{Theorem}[section]
\newtheorem{lemma}[theorem]{Lemma}
\newtheorem{proposition}[theorem]{Proposition}
\newtheorem{corollary}[theorem]{Corollary}
\theoremstyle{definition}
\newtheorem{definition}[theorem]{Definition}
\newtheorem{assumption}[theorem]{Assumption}
\theoremstyle{remark}
\newtheorem{remark}[theorem]{Remark}
\newcommand{\R}{\mathbb R}
\newcommand{\N}{\mathbb N}
\newcommand{\dd}{\mathrm d}
\DeclareMathOperator{\vol}{vol}
\DeclareMathOperator{\supp}{supp}
\DeclareMathOperator{\Area}{Area}
\DeclareMathOperator{\Scal}{Scal}
\DeclareMathOperator{\Ric}{Ric}
\DeclareMathOperator{\Hess}{Hess}
\DeclareMathOperator{\Div}{div}
\DeclareMathOperator{\tr}{tr}
\DeclareMathOperator{\sym}{sym}
\DeclareMathOperator{\id}{id}
\DeclareMathOperator{\sgn}{sgn}
\setlist{topsep=5pt,itemsep=3pt,parsep=0pt}
\setlength{\emergencystretch}{1.5em}
\allowdisplaybreaks[2]
\setcounter{tocdepth}{2}

\usepackage{accsupp}
\NewCommandCopy{\OriginalMapsto}{\mapsto}
\NewCommandCopy{\OriginalLongmapsto}{\longmapsto}
\NewCommandCopy{\OriginalHookrightarrow}{\hookrightarrow}
\renewcommand{\mapsto}{\BeginAccSupp{method=hex,unicode,ActualText=21A6}\OriginalMapsto\EndAccSupp{}}
\renewcommand{\longmapsto}{\BeginAccSupp{method=hex,unicode,ActualText=27FC}\OriginalLongmapsto\EndAccSupp{}}
\renewcommand{\hookrightarrow}{\BeginAccSupp{method=hex,unicode,ActualText=21AA}\OriginalHookrightarrow\EndAccSupp{}}
